\documentclass[11pt,a4paper]{article}
\usepackage[utf8]{inputenc}
\usepackage{amsmath, amssymb, amsfonts, bm}
\usepackage{geometry}
\geometry{margin=1in}

\title{Mathematical Foundations of the Hyperk\"ahler Quantum Stack}
\author{Advanced Computational Architecture Group}
\date{\today}

\begin{document}

\maketitle

\section{Introduction and Algebraic Core}
The framework is built upon quaternionic algebra mapping onto a compact Hyperk\"ahler manifold. A quaternion $q \in \mathbb{H}$ is expressed as:
\begin{equation}
q = w + xi + yj + zk
\end{equation}
where the imaginary units satisfy the classic Hamilton relations $i^2 = j^2 = k^2 = ijk = -1$. The non-commutative multiplication of two quaternions $q_1$ and $q_2$ preserves spatial orientation during transformation sequences. 

By definition, a Hyperk\"ahler manifold is a Riemannian manifold with three covariant constant orthogonal automorphisms $I$, $J$, and $K$ of the tangent bundle satisfying the quaternionic identities (Hitchin, 1992):
\begin{equation}
I^2 = J^2 = K^2 = IJK = -1
\end{equation}
This geometric representation serves as the primary data structure across all sequential processing steps, enforcing strict topological and rotational constraints before structural configurations are projected into lower-dimensional representations.

\section{Experiment 1: Quantum-Guided Wave Function Collapse}
This experiment bridges discrete constraint satisfaction algorithms with quadratic unconstrained binary optimization (QUBO) for micro-structural synthesis.

\subsection{Latent Encoding to Binary Quadratic Models}
Let $q_i, q_j \in \mathbb{H}$ be unit quaternions representing localized structural orientation states. The interaction weight $J_{ij}$ between distinct spatial cells is governed by the quaternionic inner product scaled by an interaction coupling constant $\alpha$:
\begin{equation}
J_{ij} = \alpha \langle q_i, q_j \rangle = \alpha (w_i w_j + x_i x_j + y_i y_j + z_i z_j)
\end{equation}
The total energy objective submitted to the quantum annealing hardware is formulated as a Binary Quadratic Model over the binary state variables $x \in \{0, 1\}^n$:
\begin{equation}
E(x) = \sum_{i=1}^n h_i x_i + \sum_{i < j} J_{ij} x_i x_j
\end{equation}
where $h_i$ represents the localized linear biases penalizing invalid boundary configurations.

\subsection{Manifold-Aware Wave Function Collapse}
The resulting optimal configuration configuration guides the multi-layer collapse operator. The continuous boundary evolution inside the grid is balanced by a stochastic propagation function parameterized by step-wise noise injection:
\begin{equation}
\mathcal{G}_{t+1} = \text{clip}\left( \mathcal{G}_t + \eta_t, -1, 1 \right)
\end{equation}
where $\mathcal{G}$ represents the multi-layer topological manifold grid and $\eta_t \sim \mathcal{N}(0, \sigma^2)$ represents localized perturbation vectors verifying stability.

\section{Experiment 2: Spatiotemporal Fluid Dynamics and Vortex Tracking}
This module maps sequential fluid properties into a geometry-aware Recurrent Neural Network (RNN) pipeline, tracking how physical fields twist and deform across a manifold over time.

\subsection{Hyperk\"ahler Long Short-Term Memory Dynamics}
The hidden state $h_t$ and cell state $c_t$ are evaluated within a quaternionic vector space. For an input vector $x_t$, the internal recurrent gates are computed using specialized linear transformations:
\begin{align}
i_t &= \sigma(W_i x_t + U_i h_{t-1} + b_i) \\
f_t &= \sigma(W_f x_t + U_f h_{t-1} + b_f) \\
o_t &= \sigma(W_o x_t + U_o h_{t-1} + b_o) \\
g_t &= \tanh(W_g x_t + U_g h_{t-1} + b_g)
\end{align}
The cell memory state update and final hidden output layer are generated via:
\begin{align}
c_t &= f_t \odot c_{t-1} + i_t \odot g_t \\
h_t &= o_t \odot \tanh(c_t)
\end{align}
where $\odot$ denotes the element-wise product, and all vector states maintain four-dimensional geometric properties mapping spatial rotations natively.

\subsection{K\"ahler Differential Field Analysis}
To track fluid vorticity, the complex-valued components undergo spatial derivative extraction. Splitting the complex state feature map into real ($r$) and imaginary ($i$) pairs, the K\"ahler differential operator computes holomorphic and anti-holomorphic derivatives using partial spatial convolution kernels $\partial_x$ and $\partial_y$:
\begin{align}
\partial z &= \frac{1}{2} \left( \frac{\partial r}{\partial x} + \frac{\partial i}{\partial y} \right) + \frac{i}{2} \left( \frac{\partial i}{\partial x} - \frac{\partial r}{\partial y} \right) \\
\bar{\partial} z &= \frac{1}{2} \left( \frac{\partial r}{\partial x} - \frac{\partial i}{\partial y} \right) + \frac{i}{2} \left( \frac{\partial i}{\partial x} + \frac{\partial r}{\partial y} \right)
\end{align}
These derivatives are evaluated across three orthogonal complex structures $I$, $J$, and $K$ to map continuous multi-dimensional fluid flow.

\section{Experiment 3: Non-Euclidean Path Planning}
This system calculates optimal trajectories over irregular or shifting operational landscapes using learnable kernel metrics.

\subsection{Symplectic Form and Learnable Hyperk\"ahler Kernels}
Given two high-dimensional spatial states $p$ and $q$, the distance function combines angular, Gaussian, and symplectic phase metrics. The canonical symplectic matrix $S$ is defined as:
\begin{equation}
S = \begin{pmatrix} 
0 & -1 & 0 & 0 \\ 
1 & 0 & 0 & 0 \\ 
0 & 0 & 0 & -1 \\ 
0 & 0 & 1 & 0 
\end{pmatrix}
\end{equation}
The localized symplectic form evaluated between states is expressed via the tensor inner product:
\begin{equation}
\omega(p, q) = p^T S q
\end{equation}
The complete learnable kernel function $K(p,q)$ driving the path trajectory evaluation is formulated as:
\begin{equation}
K(p, q) = \left[ \text{slerp\_cos}(p, q) \right]^\alpha \cdot \exp\left( -\frac{\|p - q\|^2}{2\sigma^2} \right) \cdot \cos\left( \frac{\omega(p, q)}{\tau} \right)
\end{equation}
where $\sigma$ acts as the Gaussian scaling width, $\tau$ controls the symplectic phase frequency, and $\alpha$ determines the angular constraint sharpness.

\subsection{Adaptive Potential Optimization}
The optimal trajectory minimizes the total loss function, combining the permutation target loss with a multi-part structural regularization penalty:
\begin{equation}
\mathcal{L}_{\text{total}} = \mathcal{L}_{\text{perm}} + \lambda_1 \mathcal{L}_{\text{drift}} + \lambda_2 \mathcal{L}_{\text{asym}} + \lambda_3 \mathcal{L}_{\text{canon}}
\end{equation}
The underlying matrix constraint penalties ensure that the learnable geometric space does not warp into invalid configurations during gradient descent updates:
\begin{align}
\mathcal{L}_{\text{asym}} &= \|S + S^T\|^2_F \\
\mathcal{L}_{\text{canon}} &= \|S^T S_0 S - S_0\|^2_F
\end{align}
where $S_0$ represents the baseline canonical matrix, ensuring the preservation of the underlying symplectic architecture throughout the path optimization loop (Kamenova \& Verbitsky, 2014).

\newpage
\begin{thebibliography}{9}

\bibitem{hitchin1992}
Hitchin, N. (1992). Hyperk\"ahler manifolds. \textit{S\'eminaire Bourbaki}, \textit{34}(137), 137-166.

\bibitem{kamenova2014}
Kamenova, L., \& Verbitsky, M. (2014). Families of Lagrangian fibrations on hyperk\"ahler manifolds. \textit{Advances in Mathematics}, \textit{260}, 401-413. https://doi.org/10.1016/j.aim.2013.10.033
\\ \textit{Cited by: 20}

\bibitem{debarre2018}
Debarre, O. (2018). Hyperk\"ahler manifolds. \textit{arXiv preprint arXiv:1810.02087}.
\\ \textit{Cited by: 47}

\bibitem{glover2014}
Glover, R., \& Sawon, J. (2014). Generalized twistor spaces for hyperk\"ahler manifolds. \textit{Journal of the London Mathematical Society}, \textit{91}(2), 321-342. https://doi.org/10.1112/jlms/jdu074
\\ \textit{Cited by: 12}

\end{thebibliography}

\end{document}
